\par\medskip
\SourceLeaf{072}\pdfbookmark[2]{54.}{article-054}
\begin{center}54.\end{center}
The relations of the motion to the ecliptic explained in the preceding articles clearly hold just as well even if any other plane is substituted for the ecliptic, provided the orientation of the orbital plane relative to that plane is known. The terms longitude and latitude must then be dropped. This gives rise to the problem of \textit{deriving the orientation of the orbital plane relative to a new plane from the known orientations of both planes relative to the ecliptic.} Let $n\AscendingNode$, $\AscendingNode\AscendingNode'$, and $n\AscendingNode'$ be arcs of the great circles that the plane of the ecliptic, the orbital plane, and the new plane project onto the celestial sphere (Fig. 2). To specify without ambiguity the inclination of the second circle to the third and the position of the ascending node, one of the two directions along the third circle must be chosen to correspond to the direction on the ecliptic that follows the order of the signs of the zodiac. In our figure let it be the direction from $n$ toward $\AscendingNode'$. Further, of the two hemispheres separated by the circle $n\AscendingNode'$, one must be regarded as corresponding to the northern hemisphere and the other to the southern. The two hemispheres are already distinguished by the rule that the one on the right when one proceeds along the circle in the order of the signs\footnote{That is, on the \textit{inner} surface of the sphere, which our figure represents.} is always regarded as northern. Thus in our figure $\AscendingNode$, $n$, and $\AscendingNode'$ are the ascending nodes of the second circle on the first, the third on the first, and the second on the third, respectively. The inclinations of the second to the first, the third to the first, and the second to the third are $180^0-n\AscendingNode\AscendingNode'$, $\AscendingNode n\AscendingNode'$, and $n\AscendingNode'\AscendingNode$, respectively. Our problem therefore depends on solving a spherical triangle in which the remaining parts are to be found from one side and its adjacent angles. We omit the usual rules given for this case in spherical trigonometry, since they are already well known. It is more convenient to use a different method, based on certain equations that one would seek in vain in our trigonometry books. We shall use the following equations frequently. Let $a$, $b$, and $c$ denote the sides of the spherical triangle and $A$, $B$, and $C$ the respective opposite angles.\HistoricalNoteDouble{Manuscript note:] If two triangles have the two sides $b$ and $c$ in common, then
\[
\begin{aligned}
\mathrm{I.}\quad&\frac{\sin\frac{1}{2}(a'-a)}{\sin\frac{1}{2}(A'-A)}=\frac{\sin b\sin\frac{1}{2}(C'+C)}{\cos\frac{1}{2}(B'-B)}=\frac{\sin c\sin\frac{1}{2}(B'+B)}{\cos\frac{1}{2}(C'-C)}\\[.5em]
\mathrm{II.}\quad&-\frac{\sin\frac{1}{2}(a'+a)}{\sin\frac{1}{2}(A'-A)}=\frac{\sin b\cos\frac{1}{2}(C'+C)}{\sin\frac{1}{2}(B'-B)}=\frac{\sin c\cos\frac{1}{2}(B'+B)}{\sin\frac{1}{2}(C'-C)}\\[.5em]
\mathrm{III.}\quad&-\frac{\sin\frac{1}{2}(a'-a)}{\sin\frac{1}{2}(A'+A)}=\frac{\sin b\sin\frac{1}{2}(C'-C)}{\cos\frac{1}{2}(B'+B)}=\frac{\sin c\sin\frac{1}{2}(B'-B)}{\cos\frac{1}{2}(C'+C)}\\[.5em]
\mathrm{IV.}\quad&\frac{\sin\frac{1}{2}(a'+a)}{\sin\frac{1}{2}(A'+A)}=\frac{\sin b\cos\frac{1}{2}(C'-C)}{\sin\frac{1}{2}(B'+B)}=\frac{\sin c\cos\frac{1}{2}(B'-B)}{\sin\frac{1}{2}(C'+C)}\cdot
\end{aligned}
\]
}
\SourceLeaf{073}
\[
\mathrm{I.}\quad\frac{\sin\frac{1}{2}(b-c)}{\sin\frac{1}{2}a{,}}=\frac{\sin\frac{1}{2}(B-C)}{\cos\frac{1}{2}A}
\]
\[
\mathrm{II.}\quad\frac{\sin\frac{1}{2}(b+c)}{\sin\frac{1}{2}a}=\frac{\cos\frac{1}{2}(B-C)}{\sin\frac{1}{2}A}
\]
\[
\mathrm{III.}\quad\frac{\cos\frac{1}{2}(b-c)}{\cos\frac{1}{2}a}=\frac{\sin\frac{1}{2}(B+C)}{\cos\frac{1}{2}A}
\]
\[
\mathrm{IV.}\quad\frac{\cos\frac{1}{2}(b+c)}{\cos\frac{1}{2}a}=\frac{\cos\frac{1}{2}(B+C)}{\sin\frac{1}{2}A}\cdot
\]
Although for brevity we must omit the proof of these propositions here, anyone can readily verify them for triangles in which neither the sides nor the angles exceed $180^0$. If, however, the spherical triangle is conceived in the greatest generality, with neither its sides nor its angles subject to any limits (which offers many substantial advantages, though it requires some preliminary explanation), cases may arise in which the sign must be changed in all the preceding equations. But the original signs are clearly restored as soon as one of the angles or one of the sides is increased or decreased by $360^0$. The signs as given may therefore always be safely retained, whether the remaining parts are to be determined from a side and its adjacent angles or from an angle and its adjacent sides. For our formulas will always yield either the values sought themselves or values differing from the true ones by $360^0$ and therefore equivalent to them. We reserve a fuller explanation of this matter for another occasion. For now, however, it can readily be verified by rigorous induction, that is, by a complete enumeration of all cases, that the rules we shall base on these formulas, both for solving our problem and on other occasions, hold generally in every case.
